\clearpage\thispagestyle{otFront}\phantomsection\label{references}\pdfbookmark[0]{Sources et édition}{bm-references}\otPageTitre{Sources et édition}{Références officielles · France}\fontsize{11}{15}\selectfont
\par\medskip\textbf{1. Programme actuellement applicable}\par Bulletin officiel du 30 juillet 2020, programme du cycle 4.\par\href{https://www.education.gouv.fr/bo/20/Hebdo31/MENE2018714A.htm}{\textcolor{otGreenDark}{Consulter la source officielle}}\par
\par\medskip\textbf{2. Attendus de fin de Troisième}\par Éduscol, attendus annuels de mathématiques, document de 11 pages.\par\href{https://eduscol.education.gouv.fr/sites/default/files/document/18-maths-3e-attendus-eduscol1114748pdf-74688.pdf}{\textcolor{otGreenDark}{Consulter la source officielle}}\par
\par\medskip\textbf{3. Repères annuels du cycle 4}\par Éduscol, progression annuelle et programmation par blocs.\par\href{https://eduscol.education.gouv.fr/sites/default/files/document/26-maths-c4-reperes-eduscol1114756pdf-74694.pdf}{\textcolor{otGreenDark}{Consulter la source officielle}}\par
\par\medskip\textbf{4. Calendrier des nouveaux programmes}\par Éduscol : nouveau programme de 2026 en Troisième à la rentrée 2028.\par\href{https://eduscol.education.gouv.fr/5736/ressources-d-accompagnement-du-programme-de-mathematiques-au-cycle-4}{\textcolor{otGreenDark}{Consulter la source officielle}}\par
\par\medskip\textbf{5. Définition du brevet à compter de 2027}\par Note du 11 septembre 2026, BO spécial n° 4 du 17 septembre 2026.\par\href{https://www.education.gouv.fr/bo/2026/Hebdo4/MENE2623228N}{\textcolor{otGreenDark}{Consulter la source officielle}}\par
\par\medskip\textbf{6. Différencier les entraînements}\par Éduscol : accompagnement en mathématiques en 4e, 3e et 2de.\par\href{https://eduscol.education.gouv.fr/6858/accompagner-les-eleves-en-mathematiques-4e-3e-et-2de}{\textcolor{otGreenDark}{Consulter la source officielle}}\par
\par\medskip\textbf{7. Validité des repères en 2026-2027}\par Éduscol : repères et attendus annuels, mise à jour de septembre 2026.\par\href{https://eduscol.education.gouv.fr/6910/reperes-annuels-de-progression-et-attendus-de-fin-d-annee-du-cp-la-troisieme}{\textcolor{otGreenDark}{Consulter la source officielle}}\par
\par\bigskip Sources consultées le 21 octobre 2026. Les attendus annuels et les
repères du cycle 4 ont guidé la correspondance finale. Les exercices sont
originaux, non reproduits des documents officiels. Les défis et les sujets
du livre ne sont pas des sujets officiels.\par\medskip
Le calendrier publié en 2026 est essentiel : programme nouveau en 5e dès
2026, en 4e dès 2027, en 3e dès 2028. Cette édition prépare la Troisième
2026-2027 et la session 2027 du DNB, série générale.

\clearpage
